\chapter[Redox-active sites (hemes, iron--sulfur, flavins, copper)]{Redox-active sites\\(hemes, iron--sulfur, flavins, copper)}
\label{ch:redox}

Many proteins carry a group that can take up or give away electrons: a heme,
an iron--sulfur cluster, a flavin or a copper ion. The standard Amber force
field does not know these groups, so they need special preparation before
ProPrep can build a simulation model. This chapter shows how to find such a
group with the \menupath{2a}{Redox Site Detector} and how to rename it for its
parameter set with the \menupath{4d}{Redox Site Preparer}. The worked example
is horse cytochrome~c, \pdb{1HRC}.

\begin{glance}
\glanceitem{What it does}{2a finds metals, cofactors and their protein
  ligands, and records the bonds between them. 4d renames and splits these
  residues so that they match a ready-made parameter set.}
\glanceitem{You need it when}{Your protein has a heme, an iron--sulfur
  cluster, a flavin (FMN, FAD), NAD(P), biopterin, a Zn(Cys)$_4$ site or
  another metal you want to model with bonds. You also need it before QM/MM
  (Chapter~\ref{ch:qmmm}).}
\glanceitem{Before you start}{Load a structure with \key{1a}
  (Chapter~\ref{ch:structure}). Nothing else is needed.}
\glanceitem{Where to find it}{\key{2a} Redox Site Detector, then \key{4d}
  Redox Site Preparer in the main menu.}
\glanceitem{You get}{From 2a, the site file \file{<PDB>_redox_sites.json}.
  From 4d, a renamed structure \file{final_transformed_structure.pdb} and a
  matching site file \file{..._redox_sites_updated.json}, ready for the
  Topology Generator (\key{5a}).}
\end{glance}

\section{What a redox site is}
\label{sec:redox-what}

A \term{redox center} is the atom or group that changes its number of
electrons, for example the iron of a heme. The protein holds it in place with
a few side chains, called \term{ligands}; a ligand that binds the metal from
above or below a flat heme ring is an \term{axial ligand}. In cytochrome~c the
iron has two axial ligands, His18 and Met80. A \term{redox site} is the
center together with its ligands and the bonds between them.

Hemes come in two kinds that matter here. A \term{c-type heme} (residue name
HEC) is fixed to the protein by two thioether bonds to cysteines. A
\term{b-type heme} (residue name HEM) is held only by its axial ligands. A heme
also has two acid side groups, the \term{propionates}, which can be charged
($-1$) or neutral.

\begin{behindscenes}
The Amber force field describes a residue by its name and its atom names.
A heme that is bonded to Cys14, Cys17, His18 and Met80 is one chemical unit,
but the PDB file splits it over five residues. ProPrep's parameter sets
therefore use a different cut: the heme, the bound side chains and the iron
become one residue (for example \texttt{HMO}), the propionates become separate
residues, and the four protein residues keep only their backbone. The 4d
tool does this cutting and renaming for you. The recipe it follows is called a
\term{transformer}.
\end{behindscenes}

So redox preparation has two parts:
\begin{enumerate}
  \item \textbf{Describe the site} (2a): which residues belong to it and which
        atoms are bonded. The transformer can only work if the right bonds are
        recorded.
  \item \textbf{Transform the site} (4d): choose the transformer and the
        oxidation state. ProPrep writes a renamed PDB file that the Topology
        Generator (Chapter~\ref{ch:topology}) can build.
\end{enumerate}

\section{Step by step: cytochrome c (1HRC)}
\label{sec:redox-1hrc}

This is the route that worked from start to finish. Start ProPrep with the
Full Menu (Chapter~\ref{ch:starting}).

\subsection*{Part A: load the protein}

\begin{steps}
  \item Type \key{1a} and load \typed{1HRC} from the RCSB, choosing the
        asymmetric unit (Chapter~\ref{ch:structure}).
  \item ProPrep shows a \prompt{Residue Numbering Warning}:
        \prompt{Chain A: First residue is ACE \#0}. This is normal for 1HRC.
        ACE is an acetyl cap on the first residue.
  \item Return to the main menu with \key{m}.
\end{steps}

\subsection*{Part B: detect the heme (2a)}

\begin{steps}
  \item Type \key{2a}, then \key{2} (\prompt{Detect redox-active sites}).
        A panel lists six phases: Configuration, Inventory Scan, Selection
        \& Grouping, Templates, Site Refinement, Review \& Export.
  \item Six configuration questions follow. Press \key{Enter} at each one to
        accept the defaults:
\begin{screen}
Metal elements to scan (comma-separated list, 'none', or 'all') (all):
Coordinating non-metal elements (comma-separated list, 'none', or 'all')
  (Type 'back' or 'b' to return to previous question) (all):
Search radius around each metal (Å) (range: 0.1-20.0, default: 4.0) (4.0):
Detect CYS SG–SG pairs within this distance as candidate disulfides (Å)
  (range: 1.0-5.0, default: 2.5) (2.5):
Non-standard residues (HEM, FAD, etc.)? ... (y):
Redox-active amino acids (TYR, TRP, PHE, MET)? ... (n):
\end{screen}
  \item ProPrep scans the structure and prints an inventory:
\begin{screen}
       Redox Center Inventory (2 centers found)
┏━━━━━━━━━━━━┳━━━━━━━━━━━━━━━┳━━━━━━━━━━━━━━━┳━━━━━━━━━━┓
┃ Center ID  ┃ Center Desc.  ┃ Nearby Atom   ┃ Distance ┃
┡━━━━━━━━━━━━╇━━━━━━━━━━━━━━━╇━━━━━━━━━━━━━━━╇━━━━━━━━━━┩
│ [1]        │ HEC A:105     │ MET A:80 SD   │ 2.3Å     │
│            │               │ HIS A:18 NE2  │ 2.5Å     │
...
│ [2]        │ ACE A:0       │ GLY A:1 N     │ 1.8Å     │
...
⚠ No SG–SG pairs were within the 2.50 Å auto-detection cutoff.
\end{screen}
        The heme (HEC~105) is found with Met80 and His18 next to the iron.
        ACE~0 is the acetyl cap, not a cofactor. You will leave it out in the
        next step.
  \item At \prompt{Specify disulfide bonds as index pairs \ldots{} or 'none'
        (none)} press \key{Enter}. Cytochrome~c has no disulfide
        bonds; its two cysteines are 8.7~\AA{} apart.
  \item The centers are numbered again for selection. Here ACE is 1 and HEC is
        2. Select only the heme:
\begin{screen}
│ 1      │ ACE A:0       │ —        │ organic cofactor          │
│ 2      │ HEC A:105     │ —        │ organometallic cofactor   │
...
Select additional centers to include (all): §ans«2»
\end{screen}
  \item At \prompt{Site assignments} press \key{Enter}. ProPrep makes one
        site per center (\prompt{Auto-assigning individual sites: 1:1}).
  \item Now you add the protein residues that hold the heme. At
        \prompt{Choose search constraint (1)} type \key{1} (distance), then
        \key{2} (\prompt{Minimum distance from any site atoms}), then \key{1}
        (fixed radius), then \typed{2.5} as the radius, and press \key{Enter}
        at \prompt{Use same atom filtering for search?}.
  \item ProPrep finds the four residues bonded to the heme:
\begin{screen}
│ [1]    │ MET A:80 │ 2 atoms within range │ SD S (2.28Å), SD S (2.32Å)   │
│ [2]    │ HIS A:18 │ 2 atoms within range │ NE2 N (2.47Å), NE2 N (2.04Å) │
│ [3]    │ CYS A:14 │ 1 atoms within range │ SG S (1.75Å)                 │
│ [4]    │ CYS A:17 │ 1 atoms within range │ SG S (1.86Å)                 │
...
Select residues to add: §ans«all»
Added 30 atoms from 4 complete residues to site.
\end{screen}
        The same atom can appear twice (SD at 2.28 and 2.32~\AA) because it
        is close to more than one heme atom. That is harmless.
  \item At \prompt{Define bonds between residues} type \key{1}. Enter the
        residue pairs by residue number:
        \typed{105-80 105-18 105-14 105-17}.
  \item ProPrep lists the heme atoms. For each bond, type the heme atom name,
        then pick the partner from a list sorted by distance:
  \begin{steps}
    \item Source \typed{FE}, targets \typed{1,11} (His18 NE2 at 2.04~\AA{} and
          Met80 SD at 2.32~\AA).
    \item Source \typed{CAB}, target \typed{SG}, then pick the SG of
          \textsf{CYS A:14} (1.75~\AA).
    \item Source \typed{CAC}, target \typed{SG}, then pick the SG of
          \textsf{CYS A:17} (1.86~\AA).
    \item Type \typed{done}.
  \end{steps}
\begin{screen}
│ Site Bonds (4)                                                │
│   HEC A:105 FE ↔ HIS A:18 NE2 (coordinate, 2.04Å)             │
│   HEC A:105 FE ↔ MET A:80 SD (coordinate, 2.32Å)              │
│   HEC A:105 CAB ↔ CYS A:14 SG (covalent, 1.75Å)               │
│   HEC A:105 CAC ↔ CYS A:17 SG (covalent, 1.86Å)               │
└───────────────────────────────────────────────────────────────┘
Site now contains 73 atoms, 4 bonds
Search for additional residues to add to this site? [y/n] (y): §ans«n»
\end{screen}
  \item At \prompt{Site Review Options} type \key{4} (Export sites), then
        \typed{1,2,3} to write all three file types. Then type \key{5}
        (Continue). ProPrep reports
        \prompt{Successfully detected 1 redox-active site(s)}.
\end{steps}

\begin{watchout}
When you pick a bond partner, type the atom \emph{name} (such as \typed{SG})
rather than a number. The target list is sorted again for every source atom,
so the same number points to a different atom each time. In our test, typing
\typed{28} for CAB made a ``bond'' to the SG of Cys17 at 8.21~\AA. ProPrep
accepted it without a warning. If this happens, type the source atom again and
choose \prompt{2 Replace one existing bond}.
\end{watchout}

\subsection*{Part C: add the bonds the transformer needs}

For a protein with only one site, 2a never tells you which bonds the
transformer needs. The four bonds above are \emph{not} enough. If you go
straight to 4d now, it says that no specific transformer fits
(\prompt{Compatible Transformers: None}). You can ask why:

\begin{steps}
  \item Type \key{4d}, then \key{2}.
  \item At \prompt{Enter transformer assignments} type \typed{why 1}. For the
        His/Met heme transformer you see:
\begin{screen}
  ✗ heme_his_met_axial_c_type: 12/18 requirements met — requirements not met
      - HEM-variant bonds (inter-residue post-transformation): 2 thioether-Cys
        + axial-His + axial-Met side-chain splices + 2 propionate carve-outs
        (total) (expected 6, got 0)
...
\end{screen}
  \item Press \key{Enter} on an empty line to cancel
        (\prompt{No assignments provided}) and go back with \key{m}.
\end{steps}

The six missing bonds are bonds \emph{inside} residues: the heme's C2A--CAA
and C3D--CAD (where the propionates will be cut off) and CA--CB of Cys14,
Cys17, His18 and Met80 (where the side chains will be cut off). Add them in the
site editor:

\begin{steps}
  \item Type \key{2a}, then \key{4} (\prompt{Edit detected sites}).
  \item Type \typed{edit 1}, then \key{4} (bonds), then \typed{add}.
  \item Answer the eight short questions one at a time, pressing \key{Enter}
        after each:
\begin{screen}
Atom 1:
  Chain: §ans«A»
  Resname: §ans«HEC»
  ResID: §ans«105»
  Atom name: §ans«C2A»

Atom 2:
  Chain: §ans«A»
  Resname: §ans«HEC»
  ResID: §ans«105»
  Atom name: §ans«CAA»

Distance between atoms: 1.56Å
...
Choice [1/2/3/4] (1): §ans«1»
Add covalent bond (1.56Å)? [y/n] (y): §ans«y»
✓ Bond added: A:HEC:105 C2A ↔ A:HEC:105 CAA
\end{screen}
  \item Repeat \typed{add} for \texttt{A HEC 105 C3D} -- \texttt{A HEC 105 CAD},
        and for CA -- CB of \texttt{A CYS 14}, \texttt{A CYS 17},
        \texttt{A HIS 18} and \texttt{A MET 80}. Bond type is always \key{1}
        (covalent).
  \item Type \typed{list} to check. You should see 2 coordinate bonds and 8
        covalent bonds; the six new ones are marked \texttt{intraresidue}.
  \item Type \typed{back}, \typed{back}, then \typed{save}. ProPrep answers
        \prompt{Changes saved to workspace}.
\end{steps}

\begin{watchout}
In version 1.23.0 the editor drops you out, and loses \emph{all} unsaved
edits, if a ResID answer is not a whole number. In our test we typed a whole
bond on one line at the \prompt{Chain} prompt. ProPrep accepted the line as a
chain name, the next answers slid one prompt down, and \typed{y} landed at
\prompt{ResID}. The screen showed
\texttt{Error handling menu option 'edit': invalid literal for int()} and
returned to the 2a menu. Give one answer per prompt, and type \typed{save}
after every few bonds.
\end{watchout}

\subsection*{Part D: transform the site (4d)}

\begin{steps}
  \item Type \key{4d}, then \key{2} (\prompt{Generate single redox
        microstate}). The table now lists a compatible transformer:
\begin{screen}
│ Site # │ Center │ Location         │ Other Residues      │ Compatible ...
│ 1      │ HEC    │ Chain A, HEC 105 │ 2×CYS, 1×HIS, 1×MET │ heme_his_me...
...
Available compatible transformers:
1. heme_his_met_axial_c_type (compatible with sites: 1)
2. no_transformation (compatible with sites: 1)
...
Enter transformer assignments: §ans«1»
\end{screen}
  \item Iron oxidation state: type \typed{oxidized} (ferric, Fe(III)).
        The spin state is set for you (\prompt{Setting 'low\_spin' for all
        sites}); ProPrep notes that ferric is a low-spin doublet and ferrous a
        closed-shell singlet.
  \item pH treatment: type \typed{1:fixed\_pH}
        (see Section~\ref{sec:redox-ph}).
  \item A-ring propionate: \typed{1:deprotonated}. D-ring propionate:
        \typed{1:deprotonated}. Above pH~5 both propionates are normally
        charged.
  \item Check the \prompt{Configuration Summary} and type \key{a} (apply). The
        prompt reads \prompt{Choice (a)}; the letters a/e/r/c mean apply,
        edit a step, restart and cancel.
  \item At \prompt{Review specific sites? (no)} you may type \key{1} to see
        which residue plays which role (B-ring Cys~14, C-ring Cys~17,
        His~18, Met~80, propionates A~106 and D~107). Answer \key{y} at
        \prompt{Is this component matching correct?}.
  \item ProPrep shows a plan of 12 operations and applies them. The final box
        shows the split:
\begin{screen}
│ A:HEC 105 (43 atoms) -> SPLIT:
│   ├─ A:HMO 105: FE, CHA, CHB, CHC, CHD, NA, C1A, ...
│   ├─ A:PRD 106: CAA->CA, CBA->CB, CGA->CG, O1A->O1, O2A->O2
│   └─ A:PRD 107: CAD->CA, CBD->CB, CGD->CG, O1D->O1, O2D->O2
│ A:MET 80 (8 atoms) -> SPLIT:
│   ├─ A:MEO 80: N, CA, C, O
│   └─ A:HMO 105: CB->CBM, CG->CGM
...
✓ Workspace validation passed
  • 1 centers validated
  • 73 atoms validated
  • 10 bonds validated
\end{screen}
\end{steps}

The heme is now \texttt{HMO} (ferric His/Met c-type heme), the propionates are
\texttt{PRD} (deprotonated propionate), and Cys14, Cys17, His18 and Met80 are
backbone-only residues \texttt{CYO}, \texttt{CYO}, \texttt{HIO} and
\texttt{MEO}. Appendix~\ref{app:residues} lists all these names.

\subsection*{Part E: build the topology (5a)}

Go on to the Topology Generator, \key{5a} then \key{2}
(Chapter~\ref{ch:topology}). ProPrep knows which parameter set the heme needs.
In our run we chose ff14SB, TIP3P and implicit solvent, and then the set
\prompt{1. Ferric His/Met-axial c-type heme (HMO; cyt c2) — CM5}. Option
\key{4} then built the topology and coordinate files (\file{.prmtop} and
\file{.rst7}) with \prompt{Errors = 0} and a total charge of $+7$.

\begin{tip}
Start ProPrep from the activated ProPrep environment
(Chapter~\ref{ch:installation}). In one of our sessions another AmberTools
installation came first on the computer's search path, and tLEaP ran from
there. The build still worked, but you want the AmberTools that came with
ProPrep.
\end{tip}

\section{Understanding the results}
\label{sec:redox-results}

\begin{tabularx}{\linewidth}{@{}lL@{}}
\toprule
\textbf{File} & \textbf{What it is} \\
\midrule
\file{1HRC_redox_sites.json} & The site definitions from 2a (export format 1).
  You can import it again later with 2a option 1. \\
\file{1HRC_redox_sites.pdb} & All site atoms in one PDB file (format 2). \\
\file{1HRC_site_1.pdb} & One PDB file per site (format 3). \\
\file{transformer_compatibility_report.txt} & Why each transformer did or did
  not fit (the same text as \typed{why 1}). \\
\file{final_transformed_structure.pdb} & The renamed structure. Use this for
  5a. \\
\file{final_transformed_structure_redox_sites_updated.json} & The site
  definitions with residue numbers that match the renamed PDB. \\
\file{enhanced_residue_mapping.txt} & Old $\to$ new residue and atom names. \\
\bottomrule
\end{tabularx}

\begin{watchout}
4d renumbers residues. In \file{final_transformed_structure.pdb} the ACE cap
became residue~1, so every number moved up by one: the heme is HMO~106, the
cysteines CYO~15 and 18, His HIO~19 and Met MEO~81. The mapping report and the
progress lines still use the \emph{old} numbers (HMO~105, PRD~106, PRD~107).
In multi-heme proteins the progress lines can even name a different heme
(\texttt{A\_HEC\_1299} for heme 1294). Always trust the final PDB and the
\file{_redox_sites_updated.json}, not the printed numbers.
\end{watchout}

\begin{watchout}
After one generation, all three 4d options are blocked for the rest of the
session (\prompt{Already generated microstate(s)}). To try other settings,
start a new session in a new project folder.
\end{watchout}

The propionates were written between crystal waters in the PDB file (HMO~106,
HOH~107, PRD~108, HOH~109, PRD~110), even though ProPrep printed
\prompt{No ID conflicts detected}. tLEaP built the topology without trouble,
but QM/MM set-up later picked the wrong residues. Chapter~\ref{ch:qmmm} shows
how to check and fix this.

\section{The Redox Site Detector in detail}
\label{sec:redox-detector}

The 2a menu has four options:

\begin{screen}
══ Redox Site Detector Menu ══
  1. ✓ Import redox sites from JSON file
  2. ✓ Detect redox-active sites
  3. ○ View detected site summaries
       ⚠ Need to detect or import sites first
  4. ○ Edit detected sites
       ⚠ Need to detect or import sites first
\end{screen}

Without a loaded structure, option 2 shows \prompt{Load a structure first};
option 1 always works.

\subsection{Configuration questions}

\begin{tabularx}{\linewidth}{@{}lp{2.6cm}L@{}}
\toprule
\textbf{Question} & \textbf{Default} & \textbf{Meaning} \\
\midrule
Metal elements & \texttt{all} & Which metals become centers, e.g. \typed{FE,CU}
  or \typed{none}. Typed symbols are not checked (\typed{XX} is accepted). \\
Non-metal elements & \texttt{all} & Which atoms count as ligands around a
  metal. \\
Search radius & 4.0~\AA{} (0.1--20) & Distance for the ``nearby atoms'' in the
  inventory. \\
Disulfide cutoff & 2.5~\AA{} (1--5) & Cys pairs with SG--SG below this become
  disulfide centers. \\
Non-standard residues & \texttt{y} & Include cofactors (HEM, HEC, FMN, SF4
  \ldots). \\
Redox amino acids & \texttt{n} & Also list Tyr, Trp, Phe and Met. \\
\bottomrule
\end{tabularx}

\medskip
On questions 2--6 you can type \typed{back} to return to the previous one.
Out-of-range numbers are refused with a message such as
\prompt{Value must be at most 20.0}.

\begin{watchout}
Detection lists every non-protein group as a possible cofactor. In our tests
the acetyl cap ACE (\pdb{1HRC}), glycerol GOL and acetate ACT
(\pdb{1OFW}) were all called \prompt{organic cofactor}. They are left-overs
from crystallisation or capping and are not redox-active. Leave them out at
\prompt{Select additional centers to include} by typing only the numbers you
want, for example \typed{2} or \typed{1-2}.
\end{watchout}

\subsection{Disulfides, center selection and grouping}

If no Cys pair is within the cutoff, ProPrep shows a table of all SG--SG
distances and asks \prompt{Specify disulfide bonds as index pairs (e.g. '1-2,
3-4') or 'none'}. With fewer than two cysteines the question is skipped.

At \prompt{Select additional centers to include} you can type \typed{all},
\typed{none} (disulfides only), a range \typed{1-10} or a list
\typed{1,3,5-8}.

At \prompt{Site assignments} you decide which centers form one site:
\typed{1:1 2:2} (one center each), \typed{1:1,2} (two centers in one site) or
\typed{1:1-4} (one big site). \key{Enter} gives one site per center. Every
center must be in a site; otherwise ProPrep says
\prompt{Centers not assigned to any site}.

\subsection{Templates: the way to see the bond requirements}
\label{sec:redox-templates}

When there is more than one site, ProPrep offers
\prompt{Use template-based automation for systematic redox sites? (n)}.
A \term{template} is your own label for a repeatable recipe: you set up the
first site, and ProPrep applies the same search and bonds to the others. If
the label is the name of a transformer, 4d later picks that transformer
automatically.

\begin{tip}
With two or more sites, answer \key{y} here and use the transformer name as
the site type. This is the only place where ProPrep prints the
\prompt{Bond Requirements} for the transformer, so you know exactly which bonds
to define. With a single site, use \typed{why 1} in 4d instead
(Section~\ref{sec:redox-1hrc}, Part~C).
\end{tip}

\begin{example}
The nine-heme cytochrome \pdb{1OFW} has 18 bis-His c-type hemes. We selected
hemes \typed{1-2}, pressed \key{Enter} for grouping, answered \key{y} to
templates and typed \typed{1-2:"heme\_bis\_his\_c\_type"}. ProPrep then set up
site~1 and suggested a residue-count search:
\begin{screen}
Caution: Distance cutoffs may not be satisfied at every site due to geometric
variations between sites. Residue count is generally more reliable for templates.
Choose search constraint (2): §ans«2»
Enter residue types to find (comma-separated, e.g., HIS,CYS,ASP) (HIS,CYS): §ans«HIS,CYS»
Enter number of each type to find ... (2): §ans«2,2»
...
Bond Requirements for heme_bis_his_c_type:
(At least one variant must be satisfied)
...
• HEC variant bonds (inter-residue post-transformation) (min: 6 bonds):
  Covalent:
    - HEC:C2A ↔ HEC:CAA
    - HEC:C3D ↔ HEC:CAD
    - CYS:CA ↔ CYS:CB
    - CYS:CA ↔ CYS:CB
    - HIS:CA ↔ HIS:CB
    - HIS:CA ↔ HIS:CB
\end{screen}
Same-residue pairs are allowed, so the pairs were
\typed{1294-1294 37-37 51-51 47-47 50-50}, then \typed{C2A}$\to$\typed{CAA},
\typed{C3D}$\to$\typed{CAD} and \typed{CA}$\to$\typed{CB} for each protein
residue. At \prompt{Apply these settings to all other heme\_bis\_his\_c\_type
sites (1 remaining)? (y)} we pressed \key{Enter}, and site~2 got the same six
bonds automatically.
\end{example}

\subsection{Finding the ligands: search options}

For each site ProPrep repeats a search loop until you stop it.

\begin{itemize}
  \item \textbf{Distance cutoff} (\key{1}). Choose where distances are measured
        from: \key{1} the geometric center of the site, \key{2} the nearest
        site atom (best for bonded ligands), or \key{3} atoms you pick.
        Then a fixed radius (default 5.0~\AA) or an adaptive radius that grows
        from 3 to 8~\AA{} until a target number of atoms is found.
  \item \textbf{Residue count} (\key{2}). Finds the $N$ closest residues of
        each type, e.g. \typed{HIS,CYS} with \typed{2,2}. Add MET for His/Met
        hemes. This is the most reliable choice for templates.
  \item \textbf{Finish} (\key{3}) keeps the site as it is.
\end{itemize}

At \prompt{Select residues to add} you can type numbers, ranges, \typed{all},
\typed{0} or \typed{none} (change the search and repeat), \typed{finish}, or
\typed{view} (3D viewer, not tried for this manual). Whole residues are always
added. Sites with only one residue skip the bond step
(\prompt{Need at least 2 residues to define inter-residue bonds}).

\subsection{Editing, viewing and exporting sites}

After detection, \prompt{Site Review Options} offers: \key{1} view a site
summary, \key{2} edit, \key{3} re-run detection on chosen sites, \key{4}
export, \key{5} continue. The same editor opens later with 2a option~4.

\begin{tabularx}{\linewidth}{@{}lL@{}}
\toprule
\textbf{Editor command} & \textbf{What it does} \\
\midrule
\typed{edit 1} & Open site 1. Inside: \key{1} centers, \key{2} residues,
  \key{3} atoms, \key{4} bonds, \key{5} metadata, plus \typed{summary},
  \typed{validate}, \typed{back}. \\
\typed{remove 2} & Delete a site (also \typed{1-5}, \typed{all}). \\
\typed{validate} & Check all sites
  (\prompt{Site validation passed - no issues found}). \\
\typed{save} & Keep the changes and leave. \\
\typed{cancel} & Leave without saving. \\
\bottomrule
\end{tabularx}

\medskip
In the bonds menu, \typed{list} shows coordinate and covalent bonds in
separate tables; \typed{add}, \typed{remove} and \typed{reclassify} change
them. \typed{cancel} may warn \prompt{You have unsaved changes} even if you
only pressed \key{Enter}; answer \key{y} to discard.

Export formats are \key{1} JSON (complete data, can be imported again),
\key{2} all sites in one PDB and \key{3} one PDB per site. Type a list such
as \typed{1,2,3}. An invalid number such as \typed{1,5} exports nothing
(\prompt{Invalid choices: 5}).

\begin{watchout}
2a option~1 (import) lists every JSON file in the project folder, including
ProPrep's own \file{proprep_session_*.json} files. Picking a session file
gives \prompt{Successfully imported 0 redox site(s)} followed by
\prompt{No sites found in JSON file}. Pick \file{<PDB>_redox_sites.json} or
\file{_redox_sites_updated.json}. With the right file you see, for example,
\prompt{Successfully imported 2 redox site(s)} and
\prompt{Restored transformer mappings for 1 site type(s)}.
\end{watchout}

\section{The Redox Site Preparer in detail}
\label{sec:redox-preparer}

The 4d menu has three options: \key{1} \prompt{Create transformer
(interactive PDB editor)}, \key{2} \prompt{Generate single redox microstate}
and \key{3} \prompt{Generate all redox microstates (batch processing)}.
Without sites it shows \prompt{Needs detected redox sites — run Redox Site
Detector first}.

\subsection{Built-in transformers}

\begin{tabularx}{\linewidth}{@{}lL@{}}
\toprule
\textbf{Transformer} & \textbf{For} \\
\midrule
\texttt{heme\_his\_met\_axial\_c\_type} & His/Met c-type heme (cytochrome c,
  c2) \\
\texttt{heme\_bis\_his\_c\_type} & Bis-His c-type heme (multiheme
  cytochromes) \\
\texttt{heme\_bis\_his\_b\_type} & Bis-His b-type heme (HEM, no thioether) \\
\texttt{heme\_cys\_axial\_b\_type} & Cys-ligated b-type heme (P450, NOS) \\
\texttt{iron\_sulfur\_4fe4s} & [4Fe-4S] cluster with four Cys \\
\texttt{flavin\_fmn}, \texttt{flavin\_fad} & FMN and FAD, five states
  each \\
\texttt{nicotinamide\_nad}, \texttt{nicotinamide\_nadp} & NAD$^+$/NADH,
  NADP$^+$/NADPH \\
\texttt{pterin\_biopterin} & Biopterin family, six states \\
\texttt{zinc\_cys4} & Structural Zn(Cys)$_4$ site \\
\texttt{disulfide} & Renames bonded cysteines to CYX \\
\texttt{no\_transformation} & Leaves the site as it is \\
\bottomrule
\end{tabularx}

\subsection{Assigning transformers}

At \prompt{Enter transformer assignments} type \typed{1:"name"} for each site
(several at once: \typed{1-3:"name", 4:"other"}). With one site the list
number alone works. \typed{why 2} explains, requirement by requirement, why a
transformer does not fit site~2. \key{Enter} on an empty line cancels.

\subsection{Oxidation states and spin}

After the assignment ProPrep asks one question per parameter. You can type a
value for all sites (\typed{oxidized}), per site (\typed{1:oxidized,
2:reduced}), or press \key{Enter} to be asked option by option
(\prompt{Sites for 'oxidized' (numbers/ranges, empty=skip)}).

\begin{itemize}\sloppy
  \item \textbf{Hemes}: \texttt{oxidized} or \texttt{reduced}. Spin is fixed
        (low spin for the His/Met and bis-His hemes).
  \item \textbf{Flavins} (\pdb{1FLV}): \texttt{oxidized},
        \texttt{semiquinone\_neutral} (FMNH$\cdot$),
        \texttt{semiquinone\_anionic}, \texttt{hydroquinone\_anionic},
        \texttt{hydroquinone\_neutral}. Spin follows from the state.
  \item \textbf{[4Fe-4S]} (\pdb{1CLF}): \texttt{oxidized} ([Fe$_4$S$_4$]$^{3+}$),
        \texttt{reduced\_hipip} ([Fe$_4$S$_4$]$^{2+}$) and
        \texttt{reduced\_fd} ([Fe$_4$S$_4$]$^{1+}$). Then a
        \term{broken-symmetry} parameter set: a choice between equivalent
        ways of pairing the iron spins. Codes are FO1--FO6, FR1--FR3 and
        FD1--FD6 for the three states; FO0/FR0/FD0 are averages over all of
        them. A code that does not match the state is refused, e.g.
        \prompt{'FD1' is not a valid option for site 1}.
\end{itemize}

\subsection{Fixed pH or constant pH, and the propionates}
\label{sec:redox-ph}

Heme transformers ask for a \prompt{Propionate/titratable-site pH treatment}:

\begin{itemize}\sloppy
  \item \typed{fixed\_pH}: each propionate is fixed as charged (PRD) or neutral
        (PRP). Works with the modern ff14SB and ff19SB force fields. Choose
        this for ordinary MD. You then answer the A-ring and D-ring questions
        (\texttt{deprotonated} or \texttt{protonated}).
  \item \typed{constant\_pH}: titratable propionates (PRN) from Amber's
        constant-pH and constant-redox libraries. They need the older ff10
        backbone. Choose this only for constant-pH MD
        (Chapter~\ref{ch:topology}); the propionate questions are then
        skipped.
\end{itemize}

In the batch route (option 3) the default is \texttt{constant\_pH} if you
press \key{Enter}.

\subsection{Several sites and oxidation states: batch microstates}
\label{sec:redox-batch}

A \term{microstate} is one combination of oxidation states across all sites.
For two hemes that can each be reduced or oxidized there are four. Option
\key{3} writes one PDB file per microstate. This is how we ran \pdb{1OFW}
(two hemes, after importing the sites JSON):

\begin{steps}\sloppy
  \item Type \key{4d}, then \key{3}. ProPrep assigned the transformers
        automatically from the template names given in 2a.
  \item At \prompt{Enter selection (required)} type \typed{all:1,2} (all sites
        in states 1 and 2). Other forms are \typed{1:1,2 2:1}, or
        \typed{single:1} and \typed{pairs:1} for systematic variations when
        there are many sites.
\begin{screen}
Microstate Encoding:
  Each site is represented by a 2-letter code:
    First letter: R=Reduced, O=Oxidized
    Second letter: L=Low spin, H=High spin, I=Intermediate spin
...
  1. RLRL
  2. RLOL
  3. OLRL
  4. OLOL
\end{screen}
  \item Choose \key{1} (all) or \key{2} (a subset, e.g. \typed{1,4}).
  \item Answer the pH treatment for each redox state. For the oxidized state
        we typed \typed{1-2:fixed\_pH}, \typed{1-2:deprotonated} and
        \typed{1:deprotonated,2:protonated}; for the reduced state we pressed
        \key{Enter} (constant pH).
\end{steps}

The result is \file{transformed_microstate_001.pdb}, \file{_002.pdb}, an
updated sites JSON for each, and \file{microstate_metadata.json}. ProPrep says
\prompt{Use this file for tLEaP input generation (option 3 in Topology
Generator menu)}.

\begin{watchout}
For sites with \texttt{no\_transformation} (for example Tyr or Phe in
\pdb{1UBQ}), the microstate code reads ``OL'' (oxidized/low spin) although
nothing was changed. The code is only a label here.
\end{watchout}

\subsection{Making your own transformer}

{\sloppy Option \key{1} opens a small text editor for a site's residues. Commands
include \typed{rename\_res}, \typed{rename\_atom}, \typed{movenew} (move atoms
into a new residue), \typed{param} and \typed{vary} (names that change with
the oxidation state), \typed{summary}, \typed{drop} (undo the last step),
\typed{help} and \typed{quit}. We tried the editing commands on \pdb{1OFW}.
We did not try \typed{save}, which writes the transformer to
\file{~/.proprep/transformers} and links it to a parameter library; it is not
described further here. Chapter~\ref{ch:forcefields} explains why a
transformer is also needed for parameters you import yourself.\par}

\section{Other proteins we tried}
\label{sec:redox-others}

\begin{tabularx}{\linewidth}{@{}lL@{}}
\toprule
\textbf{Protein} & \textbf{What happened} \\
\midrule
\pdb{1OFW} nine-heme cytochrome &
  18 HEC found (plus 4 GOL and 1 ACT, deselected). Two hemes done with a
  template; JSON export and import worked; batch microstates written as above.
  The \texttt{heme\_bis\_his\_c\_type} transformer gave HCO (oxidized) and
  HCR (reduced) hemes. \\
\pdb{1FLV} flavodoxin &
  One FMN found. \texttt{flavin\_fmn} with \texttt{semiquinone\_neutral}
  renamed FMN to FMS. \\
\pdb{1CLF} 2[4Fe-4S] ferredoxin &
  Two SF4 clusters, each with four Cys SG nearby. Template
  \texttt{iron\_sulfur\_4fe4s} with a count search (\typed{CYS}, \typed{4})
  and four Fe--SG bonds. One cluster became FR1 with Cys named C25--C28, the
  other FD0 with Cys named CXD. \\
\pdb{1PLC} plastocyanin (Cu) &
  The Cu and its four ligands (His37, His87, Cys84, Met92) were found with a
  3.0~\AA{} search. No transformer exists for copper; only
  \texttt{no\_transformation} is offered. Copper parameters must come from
  elsewhere (Chapter~\ref{ch:forcefields}). \\
\pdb{4HHB} hemoglobin &
  Four HEM, each with its proximal His. Grouped as one site (\typed{1:1-4}),
  no transformer fits: \prompt{Required center count (expected 1, got 4)}.
  Keep one heme per site. \\
\pdb{1UBQ} ubiquitin &
  \prompt{Found 0 potential redox centers} with the defaults. With redox
  amino acids switched on, Met1, Phe4, Phe45 and Tyr59 are listed;
  \texttt{no\_transformation} is the only choice. \\
\bottomrule
\end{tabularx}

\section{If something goes wrong}
\label{sec:redox-problems}

\begin{itemize}\sloppy
  \item \textbf{4d says \prompt{Compatible Transformers: None}.} The site lacks
    residues or bonds. Type \typed{why 1} at the assignment prompt and add
    what is missing in 2a $\to$ \key{4} $\to$ \typed{edit 1} $\to$ bonds
    (Section~\ref{sec:redox-1hrc}, Part C). For several sites, use templates
    (Section~\ref{sec:redox-templates}).
  \item \textbf{A bond of several \AA{} was created.} You typed a list number that
    pointed to another atom. Type the source atom again and choose
    \prompt{Replace one existing bond}.
  \item \textbf{The editor jumped back to the 2a menu with \texttt{invalid literal for
    int()}.} A ResID answer was not a number. Your unsaved edits are lost;
    redo them one answer per prompt.
  \item \textbf{ACE, GOL, ACT or EDO appear as cofactors.} Deselect them at center
    selection.
  \item \textbf{Residue numbers in the messages do not match the PDB.} Use the final
    PDB and the updated sites JSON file
    (Section~\ref{sec:redox-results}).
  \item \textbf{4d options are greyed out with \prompt{Already generated
    microstate(s)}.} Start a new session in a new project folder.
  \item \textbf{A menu hint reads \prompt{or press to return to the main menu}.} The
    key letter is missing from the hint in version 1.23.0. Type \key{m} for
    the main menu.
\end{itemize}

See also Chapter~\ref{ch:troubleshooting} and Chapter~\ref{ch:knownissues}.
